\chapter{Literature Review}
\label{chap:literature_review}
\label{relatedwork}
The stabilizing role of the lumbar multifidus muscle in chronic LBP has been studied extensively by Freeman et al. \cite{Freeman2010}. Their review explains how LMM dysfunction, fatty infiltration, and atrophy are related to pain and impaired spinal control. It also highlights the value of imaging techniques such as MRI and ultrasound for measuring muscle impairment and recovery.

Ardhianto \cite{DeepLearningChallenges2021} reviewed deep learning applications in muscle ultrasound imaging and discussed common difficulties such as low image quality, motion artifacts, and limited annotated data. The review grouped existing methods into segmentation and classification tasks and described architectural choices such as skip connections and improved upsampling. Although it does not focus only on LMM, the discussed methods are relevant for ultrasound-based muscle analysis.

Kim et al. \cite{Kim2020} proposed a computerized method for measuring LMM thickness and intramuscular fat from ultrasound images. Their method used contrast enhancement and Fuzzy C-Means clustering to quantify fat, reducing operator dependence. The reported measurements showed close agreement with expert measurements, and the fat quantification step provided more clinical information than ordinary histogram analysis.

The LUMINOUS database \cite{Belasso2020} is important for this project because it provides annotated lumbar multifidus ultrasound images at the L5 level. It includes hand-segmented binary masks, making it suitable for training and validating ROI detection or segmentation methods. Its limitation is that the subjects are young athletic adults, so the distribution differs from older or clinical LBP populations.

Weber et al. \cite{Weber2021} used a CNN model to segment cervical spine muscles and estimate muscle fat infiltration from Dixon MRI scans. Even though the modality is MRI rather than ultrasound, the work shows how deep learning can reduce manual effort in muscle assessment.

Rodriguez et al. \cite{Rodriguez2017} studied the effect of preprocessing on ultrasound lesion segmentation and classification. Their results showed that contrast enhancement and despeckling can strongly affect ultrasound model performance. This is relevant because the present project also observed that preprocessing choices influenced Heckmatt classification performance.

Noise handling is a separate concern in ultrasound because speckle is not just ordinary camera noise. Yu and Acton \cite{YuActon2002} proposed speckle reducing anisotropic diffusion for ultrasonic and radar images, and Loizou et al. \cite{Loizou2005} compared despeckle filters for medical ultrasound images. These works are relevant here because the images used in this project vary in contrast, grain, and machine-dependent appearance. Any preprocessing step therefore has to reduce noise without washing out the texture that is needed for Heckmatt grading.

The classification models used in the project also follow common deep learning practice. ResNet introduced residual connections to make deeper convolutional networks easier to optimize \cite{he2016resnet}, which is why ResNet18 was kept as a simple but useful baseline. The class distribution in this project was not balanced, so the work also connects with imbalanced learning studies. He and Garcia \cite{HeGarcia2009} discuss the general difficulty of learning from imbalanced data, while focal loss \cite{LinFocal2017} and class-balanced loss \cite{Cui2019} show different ways of reducing bias toward majority classes. Since medical image datasets are often small, augmentation is also important; Shorten and Khoshgoftaar \cite{Shorten2019} review image augmentation methods, and AutoAugment \cite{Cubuk2019} shows how augmentation policies can be learned from data.

Recent ultrasound foundation models have also become relevant. USFM \cite{JIAO2024103202} was trained on more than two million ultrasound images and uses a spatial-frequency masked image modelling strategy. UltraSam \cite{meyer2024ultrasam}, an ultrasound-adapted version of SAM, was trained on a large ultrasound segmentation collection and showed improved performance over general-purpose SAM on ultrasound tasks. These works suggest that ultrasound-specific pretraining may be useful for future versions of this project.

